\documentclass[a4paper,landscape,8pt]{extarticle}
\usepackage[landscape,margin=0.7cm,top=0.5cm,bottom=0.7cm]{geometry}
\usepackage[utf8]{inputenc}
\usepackage[T1]{fontenc}
\usepackage[ngerman]{babel}
\usepackage{helvet}
\renewcommand{\familydefault}{\sfdefault}
\usepackage{xcolor}
\usepackage{tikz}
\usetikzlibrary{positioning, shapes.geometric}
\usepackage{enumitem}
\usepackage{multicol}
\usepackage{array}
\usepackage{fancyhdr}
\usepackage{amssymb}

% Farben
\definecolor{phishingred}{HTML}{FF3B30}
\definecolor{spamorange}{HTML}{FF9500}
\definecolor{scampurple}{HTML}{AF52DE}
\definecolor{safegreen}{HTML}{34C759}
\definecolor{warnred}{HTML}{C62828}
\definecolor{lightgray}{HTML}{F5F5F5}
\definecolor{bordergray}{HTML}{BBBBBB}
\definecolor{appleblue}{HTML}{007AFF}

\pagestyle{fancy}
\fancyhf{}
\renewcommand{\headrulewidth}{0pt}
\fancyfoot[C]{\footnotesize Seite \thepage{} von 3}

\setlength{\parindent}{0pt}
\setlength{\parskip}{2pt}

% Kompakte Listen
\setlist{nosep, leftmargin=1em, itemsep=0pt, topsep=1pt}

\begin{document}

% ============================================================================
% SEITE 1: INFORMATIONSBLATT
% ============================================================================

\noindent{\Large\bfseries Gefahren erkennen: Phishing, Spam \& Betrug} \hfill
{\footnotesize\textbf{Lernziele:} Betrugsversuche erkennen | Richtig reagieren | Sich schützen}

\vspace{1mm}
\hrule
\vspace{1mm}

% Drei-Spalten-Layout
\begin{minipage}[t]{0.32\linewidth}
\begin{center}
\colorbox{phishingred!10}{%
\begin{minipage}{0.95\linewidth}
\vspace{2mm}
\begin{center}
{\Large\bfseries\textcolor{phishingred}{Phishing}}\\[1pt]
{\footnotesize\itshape Die gefälschte Nachricht}
\end{center}
\vspace{2mm}
\end{minipage}%
}
\end{center}

\vspace{1mm}

\textbf{Was ist Phishing?} Betrüger geben sich als \textbf{vertrauenswürdige Firma} aus (Bank, Amazon, DHL...).

\textbf{Analogie:} Wie ein \textbf{verkleideter Dieb}, der sich als Postbote ausgibt.

\textbf{Ziel der Betrüger:}
\begin{itemize}
\item Deine \textbf{Zugangsdaten} stehlen
\item Deine \textbf{Kreditkartendaten} abgreifen
\item Dich auf \textbf{gefälschte Webseiten} locken
\end{itemize}

\textbf{Typische Phishing-Nachrichten:}
\begin{itemize}
\item \glqq Ihr Konto wurde gesperrt\grqq{}
\item \glqq Bestätigen Sie Ihre Daten\grqq{}
\item \glqq Ihr Paket wartet\grqq{}
\item \glqq Sie haben gewonnen!\grqq{}
\item \glqq Letzte Mahnung!\grqq{}
\end{itemize}

\vspace{2mm}

\begin{center}
\fcolorbox{phishingred}{phishingred!8}{%
\begin{minipage}{0.92\linewidth}
\vspace{2mm}
\textbf{Goldene Regel:}\\[2pt]
Banken und seriöse Firmen fragen \textbf{niemals} per E-Mail nach Passwörtern oder PINs!
\vspace{2mm}
\end{minipage}}
\end{center}

\end{minipage}%
\hfill%
\begin{minipage}[t]{0.32\linewidth}
\begin{center}
\colorbox{spamorange!10}{%
\begin{minipage}{0.95\linewidth}
\vspace{2mm}
\begin{center}
{\Large\bfseries\textcolor{spamorange}{Spam}}\\[1pt]
{\footnotesize\itshape Unerwünschte Nachrichten}
\end{center}
\vspace{2mm}
\end{minipage}%
}
\end{center}

\vspace{1mm}

\textbf{Was ist Spam?} Massenhaft verschickte, \textbf{unerwünschte Nachrichten} -- meist Werbung oder Betrug.

\textbf{Analogie:} Wie \textbf{Werbeprospekte} im Briefkasten -- nur viel mehr und oft gefährlich.

\textbf{Arten von Spam:}
\begin{itemize}
\item \textbf{Werbe-Spam} -- nervige Angebote
\item \textbf{Betrugs-Spam} -- falsche Gewinnspiele
\item \textbf{Phishing-Spam} -- gefälschte Warnungen
\item \textbf{Viren-Spam} -- gefährliche Anhänge
\end{itemize}

\textbf{Woher haben die meine Adresse?}
\begin{itemize}
\item Datenlecks bei Firmen
\item Verkaufte Adresslisten
\item Gewinnspiele (unseriöse)
\item Zufällig generiert
\end{itemize}

\textbf{Was tun mit Spam?}
\begin{itemize}
\item \textbf{Nicht öffnen} wenn möglich
\item \textbf{Nicht antworten} -- bestätigt deine Adresse
\item \textbf{Löschen} oder als Spam markieren
\item \textbf{Keine Links} anklicken
\item \textbf{Keine Anhänge} öffnen
\end{itemize}

\end{minipage}%
\hfill%
\begin{minipage}[t]{0.32\linewidth}
\begin{center}
\colorbox{scampurple!10}{%
\begin{minipage}{0.95\linewidth}
\vspace{2mm}
\begin{center}
{\Large\bfseries\textcolor{scampurple}{Betrugsmaschen}}\\[1pt]
{\footnotesize\itshape Der Enkeltrick und mehr}
\end{center}
\vspace{2mm}
\end{minipage}%
}
\end{center}

\vspace{1mm}

\textbf{Enkeltrick (auch per WhatsApp):}\\
\glqq Hallo Oma, ich hab eine neue Nummer. Kannst du mir schnell Geld überweisen?\grqq{}

$\rightarrow$ \textbf{Immer anrufen} und nachfragen!

\vspace{2mm}

\textbf{Falsche Support-Anrufe:}\\
\glqq Hier ist Microsoft. Ihr Computer hat einen Virus.\grqq{}

$\rightarrow$ \textbf{Auflegen!} Microsoft ruft nie an.

\vspace{2mm}

\textbf{Gewinnbenachrichtigungen:}\\
\glqq Sie haben 500.000€ gewonnen! Zahlen Sie nur 50€ Gebühr.\grqq{}

$\rightarrow$ \textbf{Betrug!} Echte Gewinne kosten nichts.

\vspace{2mm}

\textbf{Liebesbetrug (Romance Scam):}\\
Online-Bekanntschaft bittet um Geld.

$\rightarrow$ \textbf{Nie Geld} an Fremde überweisen!

\vspace{2mm}

\begin{center}
\fcolorbox{scampurple}{scampurple!8}{%
\begin{minipage}{0.92\linewidth}
\vspace{2mm}
\textbf{Wenn es zu gut klingt,}\\
um wahr zu sein -- ist es meist \textbf{Betrug}.
\vspace{2mm}
\end{minipage}}
\end{center}

\end{minipage}

\vspace{1mm}

% Zusammenfassung
\begin{center}
\fcolorbox{bordergray}{lightgray}{%
\begin{minipage}{0.98\linewidth}
\centering
\vspace{1.5mm}
\textbf{Merke:} Bei Druck, Drohung oder zu guten Angeboten: \textbf{Erst denken, dann klicken}. Im Zweifel: \textbf{Jemanden fragen!}
\vspace{1.5mm}
\end{minipage}%
}
\end{center}

\newpage

% ============================================================================
% SEITE 2: WARNZEICHEN & SCHUTZ
% ============================================================================

\begin{center}
{\LARGE\bfseries Warnzeichen erkennen \& richtig reagieren}\\[3pt]
{\normalsize So entlarvst du Betrugsversuche}
\end{center}

\vspace{1mm}
\hrule
\vspace{2mm}

\begin{multicols}{2}

\section*{\textcolor{phishingred}{Warnzeichen in E-Mails}}

\textbf{So erkennst du Phishing:}

\renewcommand{\arraystretch}{1.3}
\begin{tabular}{@{}p{0.3cm}p{7cm}@{}}
$\bullet$ & \textbf{Absender prüfen:} Komische Adresse? \\
& (z.B. \texttt{sparkasse@xyz123.ru}) \\
$\bullet$ & \textbf{Anrede fehlt:} \glqq Sehr geehrter Kunde\grqq{} statt Name \\
$\bullet$ & \textbf{Rechtschreibfehler:} Seriöse Firmen prüfen das \\
$\bullet$ & \textbf{Zeitdruck:} \glqq Innerhalb 24 Stunden!\grqq{} \\
$\bullet$ & \textbf{Drohungen:} \glqq Sonst wird Ihr Konto gesperrt\grqq{} \\
$\bullet$ & \textbf{Link prüfen:} Finger drauf halten (nicht tippen!) \\
& Zeigt die echte Adresse an \\
\end{tabular}
\renewcommand{\arraystretch}{1}

\vspace{3mm}

\section*{\textcolor{warnred}{Verdächtige Links prüfen}}

\textbf{Auf dem iPhone:}
\begin{enumerate}
\item \textbf{Nicht tippen!}
\item Link \textbf{lange gedrückt halten}
\item Es erscheint eine \textbf{Vorschau} mit der echten Adresse
\item Adresse genau prüfen
\end{enumerate}

\textbf{Woran erkenne ich gefälschte Links?}
\begin{itemize}
\item \texttt{sparkasse-\textcolor{phishingred}{sicherheit}.de} -- falsch!
\item \texttt{amazon.\textcolor{phishingred}{xyz}.com} -- falsch!
\item \texttt{dhl-\textcolor{phishingred}{paket-info}.com} -- falsch!
\end{itemize}

\textbf{Echte Adressen:}
\begin{itemize}
\item \texttt{sparkasse.de}
\item \texttt{amazon.de}
\item \texttt{dhl.de}
\end{itemize}

\vspace{2mm}

\begin{center}
\fcolorbox{appleblue}{appleblue!8}{%
\begin{minipage}{0.9\linewidth}
\vspace{2mm}
\textbf{Tipp:}\\[2pt]
Gib wichtige Adressen (Bank!) immer \textbf{selbst ein} -- nie auf Links in E-Mails klicken.
\vspace{2mm}
\end{minipage}}
\end{center}

\columnbreak

\section*{\textcolor{safegreen}{Richtig reagieren}}

\textbf{Bei verdächtigen E-Mails:}
\begin{enumerate}
\item \textbf{Nicht klicken}, nicht antworten
\item Im Zweifel: \textbf{Firma direkt kontaktieren}\\
      (Telefonnummer selbst heraussuchen!)
\item E-Mail \textbf{als Spam markieren} oder löschen
\end{enumerate}

\textbf{Bei verdächtigen Anrufen:}
\begin{enumerate}
\item \textbf{Auflegen} -- das ist nicht unhöflich!
\item \textbf{Keine Daten} am Telefon nennen
\item Nummer \textbf{blockieren}
\end{enumerate}

\textbf{Bei verdächtigen WhatsApp-Nachrichten:}
\begin{enumerate}
\item Person über \textbf{alte Nummer anrufen}
\item \textbf{Persönliche Frage} stellen (die nur sie weiß)
\item Im Zweifel: \textbf{Kein Geld} überweisen
\end{enumerate}

\vspace{3mm}

\section*{Wenn du Opfer geworden bist}

\textbf{Sofort handeln:}
\begin{enumerate}
\item \textbf{Passwörter ändern} (betroffene Konten)
\item \textbf{Bank informieren} (bei Geldverlust)
\item \textbf{Anzeige erstatten} (Polizei)
\item Vertrauensperson \textbf{informieren}
\end{enumerate}

\vspace{2mm}

\begin{center}
\fcolorbox{safegreen}{safegreen!8}{%
\begin{minipage}{0.9\linewidth}
\vspace{2mm}
\textbf{Keine Scham!}\\[2pt]
Betrüger sind Profis. Es kann \textbf{jedem} passieren. Wichtig ist, \textbf{schnell zu handeln}.
\vspace{2mm}
\end{minipage}}
\end{center}

\end{multicols}

\newpage

% ============================================================================
% SEITE 3: ÜBUNGEN & CHECKLISTE
% ============================================================================

\begin{center}
{\LARGE\bfseries Übungen \& Checkliste}\\[3pt]
{\normalsize Erkenne die Gefahren}
\end{center}

\vspace{1mm}
\hrule
\vspace{2mm}

\begin{multicols}{2}

\textbf{\large Teil A: Echt oder Betrug?}

Bewerte diese Nachrichten: E = echt, B = Betrug

\vspace{2mm}

\renewcommand{\arraystretch}{1.5}
\begin{tabular}{@{}p{7.5cm}cc@{}}
& E & B \\
1. \glqq Ihr Amazon-Konto wurde gesperrt. Klicken Sie hier.\grqq{} & $\square$ & $\square$ \\
2. \glqq Hallo Oma, neue Nummer, brauch schnell Geld!\grqq{} & $\square$ & $\square$ \\
3. \glqq Hier ist Microsoft, Ihr PC hat einen Virus.\grqq{} & $\square$ & $\square$ \\
4. \glqq Sie haben 1 Mio. Euro gewonnen! Nur 50€ Gebühr.\grqq{} & $\square$ & $\square$ \\
5. \glqq Ihre Bestellung wurde versandt.\grqq{} (von DHL) & $\square$ & $\square$ \\
\end{tabular}
\renewcommand{\arraystretch}{1}

\textit{\footnotesize Lösungen: B, B, B, B, E (aber Link prüfen!)}

\vspace{4mm}

\textbf{\large Teil B: Links prüfen}

Welche Links sind verdächtig? V = verdächtig, OK = sieht gut aus

\vspace{2mm}

\renewcommand{\arraystretch}{1.5}
\begin{tabular}{@{}p{5.5cm}cc@{}}
& V & OK \\
\texttt{www.sparkasse.de} & $\square$ & $\square$ \\
\texttt{www.sparkasse-sicher.xyz} & $\square$ & $\square$ \\
\texttt{www.amazon.de} & $\square$ & $\square$ \\
\texttt{www.arnazon.de} & $\square$ & $\square$ \\
\texttt{www.dhl.de/sendung} & $\square$ & $\square$ \\
\texttt{www.dhl-paket.info} & $\square$ & $\square$ \\
\end{tabular}
\renewcommand{\arraystretch}{1}

\textit{\footnotesize Lösungen: OK, V, OK, V (m statt rn!), OK, V}

\vspace{4mm}

\textbf{\large Teil C: Was machst du?}

\vspace{2mm}

\textbf{Situation 1:} Du bekommst eine E-Mail: \glqq Ihr PayPal-Konto wurde eingeschränkt. Bestätigen Sie Ihre Daten.\grqq{}

\rule{7cm}{0.4pt}

\vspace{3mm}

\textbf{Situation 2:} Deine Enkelin schreibt per WhatsApp von neuer Nummer und bittet um 500€.

\rule{7cm}{0.4pt}

\columnbreak

\textbf{\large Teil D: Warnzeichen}

Welche Warnzeichen gibt es? Kreuze an:

\vspace{2mm}

\begin{itemize}[label=$\square$, itemsep=3pt]
\item Unpersönliche Anrede (\glqq Sehr geehrter Kunde\grqq{})
\item Rechtschreibfehler
\item Zeitdruck (\glqq Innerhalb 24 Stunden!\grqq{})
\item Drohungen (\glqq Konto wird gesperrt\grqq{})
\item Komische Absender-Adresse
\item Link führt zu anderer Adresse
\item Unerwartete Nachricht
\item Frage nach Passwort oder PIN
\end{itemize}

\textit{\footnotesize Alle sind Warnzeichen!}

\vspace{4mm}

\textbf{\large Checkliste: Das weiß ich jetzt}

\vspace{2mm}

\begin{center}
\fcolorbox{phishingred}{phishingred!10}{%
\begin{minipage}{0.95\linewidth}
\vspace{2mm}
\textbf{\textcolor{phishingred}{Phishing \& Spam}}
\begin{itemize}[label=$\square$, itemsep=2pt]
\item Ich weiß, was \textbf{Phishing} ist
\item Ich kann \textbf{Warnzeichen} erkennen
\item Ich weiß: \textbf{Nie} Passwörter per E-Mail bestätigen
\end{itemize}
\vspace{2mm}
\end{minipage}}
\end{center}

\vspace{2mm}

\begin{center}
\fcolorbox{scampurple}{scampurple!10}{%
\begin{minipage}{0.95\linewidth}
\vspace{2mm}
\textbf{\textcolor{scampurple}{Betrugsmaschen}}
\begin{itemize}[label=$\square$, itemsep=2pt]
\item Ich kenne den \textbf{Enkeltrick}
\item Ich weiß: Microsoft ruft \textbf{nie} an
\item Ich weiß: Echte Gewinne kosten \textbf{nichts}
\end{itemize}
\vspace{2mm}
\end{minipage}}
\end{center}

\vspace{2mm}

\begin{center}
\fcolorbox{safegreen}{safegreen!10}{%
\begin{minipage}{0.95\linewidth}
\vspace{2mm}
\textbf{\textcolor{safegreen}{Schutz}}
\begin{itemize}[label=$\square$, itemsep=2pt]
\item Ich kann \textbf{Links prüfen} (lange drücken)
\item Ich weiß, wie ich \textbf{richtig reagiere}
\item Ich habe eine \textbf{Vertrauensperson}, die ich fragen kann
\end{itemize}
\vspace{2mm}
\end{minipage}}
\end{center}

\end{multicols}

\vfill

\begin{center}
\footnotesize\textit{Stand: \today{} -- Im Zweifel: Nicht klicken, nicht antworten, jemanden fragen!}
\end{center}

\end{document}
